\section{Experiments}
\label{sec:experiments}

\subsection{Setup}
We run five systems ourselves under the locked protocol of \S\ref{sec:protocol} and report partner runs made under the same protocol (Table~\ref{tab:main}). Claude and GPT models run in their developers' native scaffolds, Claude Code and Codex, with web tools disabled; open-weight models and Gemini run in mini-swe-agent \citep{minisweagent2025}. Each task gets one attempt with a 50-minute budget, and the patch present at the deadline is graded in a fresh container. When the in-sandbox and fresh grades disagree, two more fresh grades are taken and the majority of the three fresh grades counts. A system is a model together with its scaffold and effort setting, so we compare systems, not models. We report 95\% Wilson intervals and paired exact McNemar tests, Holm-corrected over the ten pairs of our five runs. Versions, flags and costs are in Appendix~\ref{app:repro}.

\begin{table}[t]
\centering\small\setlength{\tabcolsep}{4pt}
\resizebox{\columnwidth}{!}{%
\begin{tabular}{llrcr}
\toprule
System & Scaffold & V2 (642) & 95\% CI & Hard (51) \\
\midrule
Claude Opus 5 (xhigh) & Claude Code & 99.4 & [98.4, 99.8] & 98.0 \\
Claude Opus 5.5 (xhigh) & Claude Code & 99.4 & [98.4, 99.8] & 96.1 \\
Claude Fable 5.1 (high) & Claude Code & 99.1 & [98.0, 99.6] & 92.2 \\
GPT-6 Astra (high) & Codex & 96.9 & [95.2, 98.0] & 90.2 \\
Gemini 3.8 Flash (high) & mini-swe-agent & 94.9 & [92.9, 96.3] & 58.8 \\
\midrule
\multicolumn{5}{l}{\emph{Partner runs (same protocol, single grade, not re-graded by us)}} \\
Kimi K3 & mini-swe-agent & 97.7 & [96.2, 98.6] & 86.3 \\
GLM-5.3 & mini-swe-agent & 95.6 & [93.8, 97.0] & 86.3 \\
Inkling & mini-swe-agent & 89.9 & [87.3, 92.0] & 31.4 \\
Claude Sonnet 5 & Claude Code & -- & -- & 88.2 \\
GPT-5.6 Terra (xhigh) & Codex & -- & -- & 86.3 \\
GPT-5.6 Sol (xhigh) & Codex & -- & -- & 82.4 \\
Claude Haiku 4.5 & Claude Code & -- & -- & 25.5 \\
\bottomrule
\end{tabular}}
\caption{Percent of tasks resolved under the locked protocol, one attempt per task, with 95\% Wilson intervals on V2. Our runs report the grade in a clean container, with disagreements settled by fresh grades.}
\label{tab:main}
\end{table}


\subsection{Main results}
\label{sec:main-results}
\paragraph{The frontier has saturated the full set.} Our five runs resolve \geminiPct--\opusXPct\% of V2 (Table~\ref{tab:main}), and every one of the \nTasks{} tasks is solved by at least one of them, which independently confirms that every task is solvable. The three Claude systems resolve \fablePct--\opusXPct\% and are statistically tied ($p\geq0.72$). Opus~5 and Opus~5.5 each fail four tasks and Fable~5.1 fails six, so any ranking among them would rest on a handful of tasks.

\paragraph{Below the frontier, V2 still separates systems.} Fable~5.1 resolves more tasks than GPT-6 Astra (\fablePct\% vs.\ \astraPct\%, Holm $p=\mcFableAstraHolm$), and the partner runs span \inklingLbPct--\kimiLbPct\% with non-overlapping intervals at the two ends. Astra and Gemini differ by two points but cannot be separated on the full set (Holm $p=\mcAstraGeminiHolm$); the Hard split separates them (\S\ref{sec:hard}). The time budget binds for some systems: \geminiTimeouts{} of Gemini's runs reached the 50-minute limit, and \geminiTimeoutsPass{} of those still passed with the patch present at the deadline.

\paragraph{The scores reflect the systems, not our edits.} Instructions edited after we inspected agent failures could in principle have been tailored to those failures. On the \nPubUnedited{} tasks whose instructions were not edited in this way, Opus~5 resolves \opusXUneditedPct\%, its full-set rate. The failures that remain are also model failures. We audited all \auditFails{} tasks that Gemini failed: the reference patch passes all of them, two extra fresh grades of each failing patch changed no grade (\auditRegrades), and Opus~5 solves \auditOpusPass{} of the \auditFails. The audit attributes \auditModel{} failures to the model, \auditBorderline{} to borderline readings of the instruction and only \auditDefects{} to residual task defects; excluding those three tasks moves Gemini from \geminiPct\% to just \geminiExDefectsPct\%.

\subsection{The Hard split}
\label{sec:hard}
A task enters the Hard split if at least two of five model families failed it in earlier runs. After removing tasks with verified defects and a final review, \nHard{} tasks remain, each passing with its reference solution and failing with an empty patch (Appendix~\ref{app:hard}).

On Hard, our runs range from \hardGeminiPct\% to \hardOpusXPct\%, and the partner runs extend the range down to \hardHaikuPct\% (Table~\ref{tab:main}). Each of the four strongest systems beats Gemini ($p<0.001$). Astra and Gemini show what the split adds: they are tied on the other 591 tasks (576 vs.\ 579, $p=0.68$), but on Hard Astra resolves \hardAstra{} tasks to Gemini's \hardGemini. Below the top four, Hard grades a broad middle tier, in which Claude Sonnet~5, GPT-5.6 Terra and Sol, Kimi~K3 and GLM-5.3 resolve 42--45 tasks, and a lower tier with Inkling (\hardInklingLb{} tasks) and Claude Haiku~4.5 (\hardHaiku). Hard does not separate the top four from each other (46--50 tasks), so ranking at the very top will need harder tasks than the public split contains. Because the selection used failures of five model families, those families' Hard scores are biased downward; systems outside the selection (Opus~5.5, Fable~5.1, Astra, Sonnet~5, Terra, Sol and Haiku~4.5) show the same spread.

\subsection{Do the data repairs matter?}
\label{sec:private}
\begin{table}[t]
\centering\small\setlength{\tabcolsep}{3pt}
\resizebox{\columnwidth}{!}{%
\begin{tabular}{lccc}
\toprule
System & As shipped & Verifier fixed & Fully fixed \\
\midrule
Opus 5 (xhigh) & 81.6 & 88.2 & \textbf{93.0} \\
Kimi K3 (max) & 78.7 & 84.2 & \textbf{88.2} \\
GLM-5.3 (max) & 77.6 & 83.1 & \textbf{87.9} \\
Gemini 3.8 Flash & 77.6 & 83.5 & \textbf{86.0} \\
Inkling (xhigh) & 67.6 & 71.3 & \textbf{79.0} \\
\bottomrule
\end{tabular}}
\caption{The never-public set (272 tasks), percent resolved. The middle column re-grades the same patches with the fixed verifier. The last column re-runs all systems after 68 instructions were corrected and the images rebuilt.}
\label{tab:private}
\end{table}

Public-split scores alone cannot show that the repairs matter, because the failure audit chose some repairs after seeing failures. We therefore applied the same fixes to a never-public set of \nPriv{} tasks from the private repositories of \nPrivOrgs{} organisations (Table~\ref{tab:private}, Appendix~\ref{app:private}). Here the instruction audit was model-blind: an LLM judge saw each instruction, its hidden tests and the reference patch, but no model output, and engineers corrected the \nPrivEdited{} instructions it flagged as unclear about what the tests check.

The repairs change scores, and they change them where they were made. Fixing the shared verifier alone raised Opus~5 from \privOpusShippedPct\% to \privOpusVPct\%, and correcting the instructions and rebuilding the images raised it to \privFinalOpusPct\%. On the \nPrivEdited{} corrected tasks, the pass rate pooled over five systems rose from \privEditedBeforePct\% to \privEditedAfterPct\%, while on the other \nPrivUnedited{} tasks it moved only from \privUneditedBeforePct\% to \privUneditedAfterPct\%. The repaired set also separates systems: Opus~5 beats every other system (Holm $p\leq0.043$), and the five systems span \privFinalInklingPct--\privFinalOpusPct\%. Re-grading missed one trick here. On two tasks that need an SDK missing from the images, three systems wrote an imitation SDK into the patch, which passed because the tests mock the client; counting those \nFakeSDK{} passes as failures lowers no score by more than 0.8 points.

\subsection{What the locked protocol catches}
\label{sec:protocol-results}
\paragraph{No retrieval.} With the network open, \netTrajOpen{} of \nTasks{} Opus~5 runs (\netTrajOpenPct\%) contacted code hosts and \netOwnRepoOpen{} downloaded the task's own repository (Appendix~\ref{app:exploits}). Under the lock, agents still tried the network, issuing between \netOpusFiveFive{} and \netAstra{} fetch commands per run, mostly package installs. We checked every flagged command by hand and audited every tool call of the ten locked runs (\relayCallsPub{} public, \relayCallsPriv{} private): all requests were refused, and no run retrieved external content. Gemini once disguised a request to the Go module proxy as traffic to the model endpoint; it was refused as well.

\paragraph{Forged dependencies.} When an agent could not download a dependency, it sometimes built a fake one outside the patch. Gemini fabricated \texttt{go.sum} and \texttt{go.work.sum} entries or edited the shared module cache in place in \nGeminiForgeries{} tasks, and in earlier runs Opus~5 and Inkling forged a module version with a fabricated \texttt{go.sum} entry (Table~\ref{tab:flips}). All of these passed in the agent's sandbox and failed in the clean container, because the fake lived outside the patch; without the second grade, each would have counted as a pass. Across our five V2 runs, the two grades disagreed on one to five tasks per run, and every disagreement other than the forgeries was a flaky test or a sandbox-side flake, settled by the fresh-grade majority.

\subsection{Effect of the agent harness}
\label{sec:harness}
\todo{Placeholder. Planned analysis: the same models under different scaffolds on the Hard split, all under the locked protocol. Preliminary team runs are in Appendix~\ref{app:harness} (Table~\ref{tab:harness}); failing runs used more tokens than passing ones in \hsFailMoreTokens{} model--harness pairs. Reconcile the mini-swe-agent cells with Table~\ref{tab:main} (Inkling: \hardInklingMini{} vs.\ \hardInklingLb{} Hard tasks) and report effort settings and cost.}
